% 作者 shingja · github.com/shingja/shingja-math-modeling-workbench · 双许可见 LICENSE · 保留版权标识
% ============================================================
% 06 — 灵敏度分析
% 说明: 检验关键参数变化对模型结果的影响程度。
%       至少包含参数扰动分析和鲁棒性检验。
%       数值必须从 results/results.json 提取, 禁止手写。
%       数值格式: 4位有效数字, 百分比2位小数, p值3位小数。
% ============================================================

\section{灵敏度分析}

为验证模型的稳定性和鲁棒性，本节对模型中的关键参数进行灵敏度分析。

\subsection{参数扰动分析}
% [此处分析关键参数在合理范围内变化时, 目标函数/决策变量的变化。]

选取模型中[参数1]、[参数2]和[参数3]作为关键参数，在基准值的 $\pm 10\%$ 范围内进行扰动分析。结果如表\ref{tab:sensitivity}所示。

\begin{table}[H]
    \centering
    \caption{关键参数灵敏度分析}
    \label{tab:sensitivity}
    \begin{tabular}{ccccc}
        \toprule
        \textbf{参数} & \textbf{变化范围} & \textbf{目标值变化} & \textbf{变化率} & \textbf{灵敏度等级} \\
        \midrule
        [参数1] & $[-10\%, +10\%]$ & [范围] & [百分比] & [低/中/高] \\
        [参数2] & $[-10\%, +10\%]$ & [范围] & [百分比] & [低/中/高] \\
        [参数3] & $[-10\%, +10\%]$ & [范围] & [百分比] & [低/中/高] \\
        \bottomrule
    \end{tabular}
\end{table}

分析结果表明，[参数1]对模型结果影响[程度]，灵敏度等级为[低/中/高]；而[参数2]和[参数3]的影响[程度]。

\subsection{鲁棒性检验}
% [此处从不同角度验证模型的可靠性: 不同初始值/不同数据子集/不同随机种子等。]

为进一步验证模型的鲁棒性，本文进行了以下检验：

\begin{enumerate}
    \item \textbf{初始值敏感性:} 使用[数量]组不同初始值重复求解，目标函数值的标准差为[数值]（相对标准差为[百分比]），表明算法对初始值[不/较/非常]敏感。
    \item \textbf{数据扰动检验:} 对输入数据添加[水平]的高斯噪声，重复[次数]次实验，结果变化在[范围]内，表明模型对数据噪声[鲁棒性程度]。
\end{enumerate}

综合灵敏度分析结果，模型在关键参数扰动下表现[稳定/较稳定/不稳定]，验证了模型的[可靠性/有效性]。
